\let\negmedspace\undefined
\let\negthickspace\undefined
\documentclass[journal]{article}
\usepackage[a5paper, margin=10mm, onecolumn]{geometry}
\usepackage{lmodern} % Ensure lmodern is loaded for pdflatex

\setlength{\headheight}{1cm} % Set the height of the header box
\setlength{\headsep}{0mm}     % Set the distance between the header box and the top of the text

\usepackage{gvv-book}
\usepackage{gvv}
\usepackage{cite}
\usepackage{textcomp}
\usepackage{amsmath,amssymb,amsfonts,amsthm}
\usepackage{algorithmic}
\usepackage{graphicx}
\graphicspath{{./figs/}}
\usepackage{textcomp}
\usepackage{xcolor}
\usepackage{txfonts}
\usepackage{listings}
\usepackage{enumitem}
\usepackage{mathtools}
\usepackage{gensymb}
\usepackage{comment}
\usepackage[breaklinks=true]{hyperref}
\usepackage{tkz-euclide} 
\usepackage{listings}
\usepackage{gvv}                                        
\def\inputGnumericTable{}                                 
\usepackage[latin1]{inputenc}                                
\usepackage{color}                                            
\usepackage{array}                                            
\usepackage{longtable}                                       
\usepackage{calc}                                             
\usepackage{multirow}                                         
\usepackage{hhline}                                           
\usepackage{ifthen}                                           
\usepackage{lscape}
\usepackage{circuitikz}
\tikzstyle{block} = [rectangle, draw, fill=blue!20, 
text width=4em, text centered, rounded corners, minimum height=3em]
\tikzstyle{sum} = [draw, fill=blue!10, circle, minimum size=1cm, node distance=1.5cm]
\tikzstyle{input} = [coordinate]
\tikzstyle{output} = [coordinate]


\begin{document}
	
	\bibliographystyle{IEEEtran}
	\vspace{3cm}
	
\title{5.9.1}
\author{EE25BTECH11047 - RAVULA SHASHANK REDDY}
\maketitle
\hrulefill
\bigskip 

\renewcommand{\thetable}{\theenumi}
\setlength{\intextsep}{10pt}

\textbf{Question:} \\

 Consider the basis 
\[
\{\vec{u}_1, \vec{u}_2, \vec{u}_3\} \text{ of } \mathbb{R}^3, \text{ where }
\vec{u}_1 = \myvec{1\\0\\0}, \ 
\vec{u}_2 = \myvec{1\\1\\0}, \
\vec{u}_3 = \myvec{1\\1\\1}.
\]
Let $\{\vec{f}_1, \vec{f}_2, \vec{f}_3\}$ be the dual basis and 
$f(a,b,c) = a+b+c$. If 
\[
f = \alpha_1 \vec{f}_1 + \alpha_2 \vec{f}_2 + \alpha_3 \vec{f}_3,
\] 
find $\myvec{\alpha_1& \alpha_2 &\alpha_3}$.\\

\textbf{Solution:}\\

The dual basis satisfies :
\begin{align}
\vec{f}_i(\vec{u}_j) &= \delta_{ij} = 
\begin{cases} 
1, & i=j \\ 
0, & i\neq j 
\end{cases}. \label{eq:delta}
\end{align}

\begin{align}
\vec{U} &= \myvec{\vec{u}_1 \ \vec{u}_2 \ \vec{u}_3} 
= \myvec{1 & 1 & 1\\ 0 & 1 & 1\\ 0 & 0 & 1}, \\
\vec{F} &= \myvec{\vec{f}_1 \\ \vec{f}_2 \\ \vec{f}_3}.
\end{align}

Then the dual basis condition becomes
\begin{align}
\vec{F} \, \vec{U} &= \vec{I_3}. \label{eq:F_U}\\
\vec{F}^{-1} &= \vec{U} 
\end{align}


\begin{align}
\vec{f} &= \alpha_1 \vec{f}_1 + \alpha_2 \vec{f}_2 + \alpha_3 \vec{f}_3 
\implies \vec{f} = \vec{a}^T \vec{F}, \\
\text{where}\quad\vec{a}&=\myvec{\alpha_1\\ \alpha_2\\\alpha_3} \quad \vec{f}=\myvec{1&1&1}  \\
\vec{a}^T \vec{F} &= \vec{f} \implies \vec{a}^T = \vec{f} \vec{F}^{-1} = f \vec{U}, \\
\vec{a}^T &= \myvec{1 & 1 & 1} \myvec{1 & 1 & 1\\ 0 & 1 & 1\\ 0 & 0 & 1} 
= \myvec{1 & 2 & 3}.
\end{align}

\begin{align}
\myvec{\alpha_1\\ \alpha_2\\\alpha_3} = \myvec{1\\2\\3}
\end{align}

\end{document}
